% ============================================================
%  CHAPTER 3 — SENTENCE PATTERNS
% ============================================================
\section{Sentence Patterns}

Each entry shows a \textbf{formula} (shaded box), one or more \textbf{examples}, and an \textbf{English gloss}.
Add new examples from sessions directly to the tables.

% Column types for pattern tables
\newcolumntype{H}{>{\raggedright\arraybackslash}p{4.0cm}}
\newcolumntype{R}{>{\raggedright\arraybackslash}p{3.4cm}}
\newcolumntype{G}{>{\raggedright\arraybackslash}p{3.6cm}}

% ------------------------------------------------------------
\subsection{3.1 Identity \& Description — \textit{honaa}}

\pattern{Subject + [Noun / Adjective] + honaa (conjugated)}

\begin{tabular}{@{}HRG@{}}
  \toprule
  Hindi & Transliteration & English \\
  \midrule
  \hw{मैं} student \hw{हूँ.}       & main student hoon      & I am a student. \\
  \hw{वह} doctor \hw{है.}          & vah daaktar hai      & He/She is a doctor. \\
  \hw{आप अच्छे हैं.}          & aap acche hain          & You are good. \\
  \hw{यह किताब अच्छी है.}     & yah kitaab acchii hai   & This book is good. \\
  \bottomrule
\end{tabular}

% ------------------------------------------------------------
\subsection{3.2 Possession}

\subsubsection*{Possession with \textit{ke paas}}

\pattern{Subject + ke paas + Object + honaa}

\begin{tabular}{@{}HRG@{}}
  \toprule
  Hindi & Transliteration & English \\
  \midrule
  \hw{मेरे पास किताब है.}     & mere paas kitaab hai    & I have a book. \\
  \hw{उसके पास गाड़ी है.}     & uske paas gaadii hai   & He/She has a car. \\
  \hw{क्या तुम्हारे पास पैसे हैं?} & kyaa tumhaare paas paise hain? & Do you have money? \\
  \bottomrule
\end{tabular}

\subsubsection*{Belonging with \textit{kaa / kii / ke}}

\pattern{Owner + kaa/kii/ke + Possessed noun}

\noindent\textit{kaa/kii/ke} agrees with the \emph{possessed} noun, not the owner.

\begin{tabular}{@{}HRG@{}}
  \toprule
  Hindi & Transliteration & English \\
  \midrule
  \hw{मेरा नाम John है.}       & meraa naam John hai     & My name is John. \\
  \hw{उसकी किताब अच्छी है.}   & uski kitaab acchii hai  & Her book is good. \\
  \hw{हमारा घर बड़ा है.}       & hamaaraa ghar badaa hai & Our house is big. \\
  \bottomrule
\end{tabular}

% ------------------------------------------------------------
\subsection{3.3 Simple Actions — Present Habitual}

\pattern{Subject + [Object] + verb-stem + taa/tii/te + hoon/hai/hain/ho}

\begin{tabular}{@{}HRG@{}}
  \toprule
  Hindi & Transliteration & English \\
  \midrule
  \hw{मैं चाय पीता हूँ.}       & main chaay piitaa hoon  & I drink tea. (m.) \\
  \hw{वह रोज़ स्कूल जाती है.}  & vah roz skuul jaatii hai & She goes to school every day. \\
  \hw{हम हिंदी सीखते हैं.}     & ham Hindi sikhte hain    & We learn Hindi. \\
  \hw{तुम क्या करते हो?}       & tum kyaa karte ho?       & What do you do? \\
  \bottomrule
\end{tabular}

% ------------------------------------------------------------
\subsection{3.4 Wants \& Desires}

\subsubsection*{Wanting/needing something — \hw{चाहिए}}

\pattern{Subject (oblique) + ko + Object + chaahiye}

\begin{tabular}{@{}HRG@{}}
  \toprule
  Hindi & Transliteration & English \\
  \midrule
  \hw{मुझे पानी चाहिए.}        & mujhe paanii chaahiye    & I want / need water. \\
  \hw{आपको क्या चाहिए?}        & aapko kyaa chaahiye?     & What do you need? \\
  \hw{उसे आराम चाहिए.}         & use aaraam chaahiye      & He/She needs rest. \\
  \bottomrule
\end{tabular}

\subsubsection*{Wanting to do something — \hw{चाहना}}

\pattern{Subject + Infinitive + chaahtaa/chaahte/chaahti + hoon/hai/hain}

\begin{tabular}{@{}HRG@{}}
  \toprule
  Hindi & Transliteration & English \\
  \midrule
  \hw{मैं जाना चाहता हूँ.}     & main jaanaa chahtaa hoon & I want to go. (m.) \\
  \hw{वह खाना चाहती है.}       & vah khaanaa chahtii hai   & She wants to eat. \\
  \hw{हम Hindi सीखना चाहते हैं.} & ham Hindi sikhnaa chahte hain & We want to learn Hindi. \\
  \bottomrule
\end{tabular}

% ------------------------------------------------------------
\subsection{3.5 Ability — \hw{सकना}}

\pattern{Subject + verb-stem + sak + taa/tii/te + hoon/hai/hain}

\begin{tabular}{@{}HRG@{}}
  \toprule
  Hindi & Transliteration & English \\
  \midrule
  \hw{मैं हिंदी बोल सकता हूँ.}  & main Hindi bol saktaa hoon & I can speak Hindi. \\
  \hw{क्या आप आ सकते हैं?}      & kyaa aap aa sakte hain?     & Can you come? \\
  \hw{वह नहीं तैर सकती.}        & vah nahin tair saktii      & She cannot swim. \\
  \bottomrule
\end{tabular}

% ------------------------------------------------------------
\subsection{3.6 Questions}

\subsubsection*{Yes/No — add \hw{क्या} at the start}

\pattern{kyaa + [statement]?}

\begin{tabular}{@{}HRG@{}}
  \toprule
  Hindi & Transliteration & English \\
  \midrule
  \hw{क्या आप ठीक हैं?}        & kyaa aap thiik hain?   & Are you well? \\
  \hw{क्या यह सच है?}          & kyaa yah sach hai?      & Is that true? \\
  \hw{क्या आप Hindi जानते हैं?} & kyaa aap Hindi jaante hain? & Do you know Hindi? \\
  \bottomrule
\end{tabular}

\subsubsection*{Information questions — question word replaces the unknown}

\pattern{[Question word] + Subject + [Object] + Verb?}

\begin{tabular}{@{}HRG@{}}
  \toprule
  Hindi & Transliteration & English \\
  \midrule
  \hw{आप कहाँ जाते हैं?}       & aap kahaan jaate hain? & Where do you go? \\
  \hw{यह क्या है?}             & yah kyaa hai?           & What is this? \\
  \hw{आप कैसे हैं?}            & aap kaise hain?         & How are you? \\
  \hw{वह कौन है?}              & vah kaun hai?           & Who is that? \\
  \hw{तुम कब आओगे?}            & tum kab aaoge?          & When will you come? \\
  \bottomrule
\end{tabular}

% ------------------------------------------------------------
\subsection{3.7 Negation}

\pattern{Subject + [Object] + nahin + Verb \ \ \textnormal{(present)}}

\begin{tabular}{@{}HRG@{}}
  \toprule
  Hindi & Transliteration & English \\
  \midrule
  \hw{मैं नहीं जाता.}          & main nahin jaataa      & I don't go. \\
  \hw{मुझे नहीं पता.}          & mujhe nahin pataa      & I don't know. \\
  \hw{वह चाय नहीं पीती.}       & vah chaay nahin piitii & She doesn't drink tea. \\
  \bottomrule
\end{tabular}

% ------------------------------------------------------------
\subsection{3.8 [Add New Pattern Here]}

\textit{[Add a title, copy the pattern and table below, and fill in examples from your sessions.]}

\pattern{Formula goes here}

\begin{tabular}{@{}HRG@{}}
  \toprule
  Hindi & Transliteration & English \\
  \midrule
  % \hw{example} & romanisation & translation \\
  \bottomrule
\end{tabular}

% ============================================================
%  SESSION LOG
% ============================================================
\clearpage
\section{Session Log}

\textit{Append one row after each lesson.}

\begin{longtable}{@{}p{2.4cm} p{3.8cm} p{8cm}@{}}
  \toprule
  Date & Topic & Notes / New items \\
  \midrule
  \endfirsthead
  \toprule
  Date & Topic & Notes \\
  \midrule
  \endhead
  \bottomrule
  \endfoot
  % example row (uncomment and edit):
  % 2025-01-07 & Greetings, \textit{honaa} & Covered namaste, present tense of honaa. \\[4pt]
\end{longtable}
